% =====================================================================
\chapter{Đóng góp của thành viên và phân chia điểm}\label{chap:donggop}
% =====================================================================

\section{Cách nhóm chia việc}
Đề tài có ba mảng khối lượng tương đương: \textbf{phần mềm} (firmware ba bo mạch,
trang web, kiểm thử), \textbf{điện} (chọn linh kiện, nguồn, đi dây) và \textbf{cơ khí}
(mô hình, cơ cấu thanh chắn). Mỗi thành viên phụ trách chính một mảng và tham gia
phần giao giữa các mảng: người làm điện viết firmware đọc cảm biến mình lắp, người làm
mô hình viết firmware điều khiển servo mình gắn. Nhóm làm việc trên một kho GitHub
chung từ 22/09 đến 30/10/2026, họp mỗi tuần để tích hợp và chạy thử trên mô hình.

\begin{table}[H]
\centering\small
\caption{Mảng phụ trách chính}\label{tab:mang}
\begin{tabularx}{\linewidth}{@{}lXX@{}}
\toprule
\textbf{Thành viên} & \textbf{Mảng chính} & \textbf{Phần giao với mảng khác} \\ \midrule
Nguyễn Chí Minh & Phần mềm: firmware máy chủ, trang web, giao thức, bảo mật, kiểm thử & Tích hợp firmware hai client vào giao thức chung; tối ưu toàn hệ thống \\
Đinh Đức Thịnh & Điện: chọn phần cứng, thiết kế nguồn, sơ đồ chân, đi dây & Firmware RFID Client và Sensor Client; khởi tạo dự án ESP-IDF \\
Nguyễn Việt Tiến & Cơ khí: mô hình bãi xe, cơ cấu thanh chắn servo & Firmware thanh chắn và còi; tài liệu hướng dẫn cài đặt \\
\bottomrule
\end{tabularx}
\end{table}

\section{Chi tiết đóng góp từng thành viên}

\subsection{Nguyễn Chí Minh (24020232) — phần mềm và tích hợp}
\begin{itemize}
  \item \textbf{Firmware máy chủ ESP32-C3}: dựng kiến trúc phân lớp, các tác vụ FreeRTOS
        và cơ chế đồng bộ; UDP server, HTTP server với 12 route; mDNS
        \code{parking.local} tự thông báo lại khi đổi IP; xử lý lỗi \code{accept (23)}
        bằng giới hạn và thu hồi kết nối.
  \item \textbf{Màn hình ST7789}: driver SPI có DMA, giao diện \code{PARKSENSE}, hóa đơn,
        thông báo quẹt sai cổng, màn hình cháy; nén ảnh động LZSS và tiền xử lý trên máy
        tính để giảm RAM và CPU.
  \item \textbf{Giao thức và bảo mật}: định dạng JSON và bộ đọc JSON không cấp phát động;
        ký gói SipHash-2-4, cửa sổ chống phát lại, ACK cho sự kiện quẹt thẻ; mã PIN cho
        API; lưu phiên vào NVS và đồng hồ bãi lấy giờ từ trình duyệt.
  \item \textbf{Trang web giám sát}: sơ đồ bãi trực quan, điều khiển thanh chắn, nhật ký
        vào/ra kèm phí, cảnh báo lệch thẻ/cảm biến.
  \item \textbf{Kiểm thử}: lớp giả lập Arduino để chạy nguyên mã hai sketch trên máy tính;
        6 bộ kiểm thử, 318 phép kiểm tra; công cụ giả lập client bằng Python.
  \item \textbf{Báo cáo}: Chương 3, 5 và các sơ đồ Mermaid.
\end{itemize}
\textit{Các commit tiêu biểu}:
\begin{itemize}[nosep]
  \item \code{init server}, \code{add udp server}, \code{add display st7789},
        \code{add kaoruko anime gif to display}, \code{add emergency button};
  \item \code{update dashboard UI and display UI, fix rgb565 byte swap logic};
  \item \code{Add parking.local mDNS}, \code{Optimize firmware, add host tests},
        \code{Fix HTTP accept failures}.
\end{itemize}

\subsection{Đinh Đức Thịnh (24020313) — điện, phần cứng và client}
\begin{itemize}
  \item \textbf{Lựa chọn phần cứng}: so sánh và chốt ESP32-C3 SuperMini cho máy chủ,
        ESP8266 cho client; cảm biến IR HW-201, module lửa, cảm biến mực nước, RC522,
        DHT11 theo tiêu chí giá, mức logic 3,3\,V và số chân (Bảng~\ref{tab:sensors}).
  \item \textbf{Thiết kế nguồn}: một điểm cấp nguồn qua jack DC trên vách sau; diode
        chống cắm ngược cực; tụ hóa lọc sụt áp khi servo khởi động — nguyên nhân làm
        ESP32-C3 tự reset ở bản lắp thử đầu tiên; tách nguồn 5\,V cho servo, nối chung GND.
  \item \textbf{Sơ đồ chân và đi dây}: lập sơ đồ chân ba bo tránh chân strapping và UART
        (Bảng~\ref{tab:pin-srv}, \ref{tab:pin-clients}); gắn 5 module IR và module lửa ở
        mặt ngoài vách để dây không đi qua lòng bãi; đổi chân GPIO client khi phát hiện
        xung đột lúc khởi động (commit \code{Change some client GPIO}); chỉnh biến trở
        từng module IR theo độ phản xạ của xe mẫu.
  \item \textbf{Firmware client}: Sensor Client đọc 5 IR, lửa, mực nước; RFID Client với
        hai đầu đọc RC522 dùng chung SPI bật/tắt ăng-ten luân phiên để tránh nhiễu chéo,
        DHT11, điểm truy cập Wi-Fi, OTA; logic vào/ra theo cổng ở máy chủ.
  \item \textbf{Báo cáo}: Chương 2, 4 và phần nguồn điện, sơ đồ chân.
\end{itemize}
\textit{Các commit tiêu biểu}: \code{Init New ESP-IDF Project},
\code{Restructure ESP32 server and add client sketch}, \code{Change some client GPIO},
\code{Update client RFID}, \code{Complete RFID in and out logic}.

\subsection{Nguyễn Việt Tiến (24020322) — mô hình và thanh chắn}
\begin{itemize}
  \item \textbf{Mô hình bãi xe} (Hình~\ref{fig:model}): thiết kế kích thước 5 ô đỗ và hai
        làn theo chiều dài xe mẫu; cắt, dán đế và vách bằng bìa formex; khoét lỗ cho đầu
        phát/thu của cảm biến đúng giữa từng ô; làm hộp điều khiển trung tâm có ô cho màn
        hình, hai hông gắn đầu đọc RC522, nóc gắn DHT11.
  \item \textbf{Cơ cấu thanh chắn}: gắn hai servo SG90 ở đầu vách, làm cần chắn nhẹ để
        servo không quá tải; căn góc đóng/mở thực tế cho hai cổng lắp đối xứng (commit
        \code{reverse door}: cổng ra quay ngược chiều cổng vào).
  \item \textbf{Firmware thanh chắn và còi}: điều khiển hai servo bằng LEDC, tự đóng sau
        khi mở bằng thẻ, chế độ mở khi cháy (commit \code{Add dual SG90 barrier control and
        test sketch}, \code{update automated door}); module còi báo cháy
        (\code{implement buzzer}).
  \item \textbf{Tài liệu và kiểm thử phần cứng}: README hướng dẫn cài đặt
        (\code{docs: add project setup README}); kịch bản kiểm thử trên mô hình
        (Bảng~\ref{tab:systest}) và quay video demo.
  \item \textbf{Báo cáo}: Chương 1, 6, 7, 8.
\end{itemize}

\subsection{Cả nhóm}
Phân tích yêu cầu và 10 use case; thống nhất giao thức giữa ba bo; các buổi tích hợp
trên mô hình (tinh chỉnh thời gian lọc nhiễu IR từ 250\,ms xuống 50\,ms, thời gian mở
cổng từ 5\,s xuống 3\,s, thứ tự ô S1–S5); rà soát chéo báo cáo.

\section{Lịch sử đóng góp trên GitHub}
\begin{table}[H]
\centering\small
\caption{Các mốc chính trong lịch sử commit}\label{tab:timeline}
\begin{tabularx}{\linewidth}{@{}llX@{}}
\toprule
\textbf{Thời gian} & \textbf{Thành viên} & \textbf{Nội dung} \\ \midrule
22–25/09 & Thịnh & Khởi tạo dự án ESP-IDF, thử nghiệm bo mạch \\
29/09–08/10 & Minh & Khung máy chủ, UDP, HTTP, màn hình ST7789, trang web \\
02/10 & Thịnh & Tổ chức lại thư mục máy chủ, sketch client đầu tiên \\
13–20/10 & Minh, Thịnh & Định dạng bản tin, logic cảm biến IR, đổi chân GPIO client \\
20–27/10 & Tiến & Hai thanh chắn servo, tự đóng, còi, README \\
27–30/10 & Thịnh & RFID Client hai đầu đọc, logic vào/ra theo cổng \\
29–30/10 & Minh & Tối ưu, mDNS, kiểm thử, bảo mật, lưu NVS, báo cáo \\
\bottomrule
\end{tabularx}
\end{table}

Số commit không phản ánh đầy đủ khối lượng công việc: phần điện (chọn linh kiện, đo
đạc, đi dây, sửa lỗi nguồn) và phần cơ khí (chế tạo mô hình, căn chỉnh servo) chiếm
nhiều thời gian nhưng không nằm trong Git. Vì vậy nhóm đánh giá đóng góp theo mảng
việc và thời gian bỏ ra, không theo số dòng mã.

\section{Tỉ lệ đóng góp và phân chia điểm}
Cả ba thành viên thống nhất rằng \textbf{mức đóng góp là ngang nhau}. Theo quy định
của môn học, toàn bộ thành viên nhận 100\% số điểm của báo cáo.

\begin{table}[H]
\centering\small
\caption{Tỉ lệ đóng góp và tỉ lệ điểm}\label{tab:diem}
\begin{tabular}{@{}llcc@{}}
\toprule
\textbf{Thành viên} & \textbf{MSSV} & \textbf{Tỉ lệ đóng góp} & \textbf{Tỉ lệ điểm báo cáo} \\ \midrule
Nguyễn Chí Minh & 24020232 & 33,3\% & 100\% \\
Đinh Đức Thịnh & 24020313 & 33,3\% & 100\% \\
Nguyễn Việt Tiến & 24020322 & 33,3\% & 100\% \\
\bottomrule
\end{tabular}
\end{table}

\vspace{1cm}
\noindent Xác nhận của các thành viên:
\par\vspace{0.5cm}
\noindent
\begin{tabular*}{\linewidth}{@{\extracolsep{\fill}}ccc@{}}
Nguyễn Chí Minh & Đinh Đức Thịnh & Nguyễn Việt Tiến \\[1.6cm]
\rule{3.6cm}{0.4pt} & \rule{3.6cm}{0.4pt} & \rule{3.6cm}{0.4pt} \\
\end{tabular*}
